\ifdefined\iscombined
	\lecturetitle{Formal Modelling}
\else
\documentclass{article}
\usepackage{changepage}
\usepackage{centernot}
\usepackage{listings}
\usepackage{amsthm}
\usepackage{braket}
\usepackage{comment}
\usepackage{algorithm}
\usepackage{algpseudocode}
\usepackage{amsmath}
\usepackage{amsfonts}
\usepackage{sectsty}
\usepackage{hyperref}
\usepackage{fancyhdr}
\usepackage[top=1in,bottom=1in,left=1in,right=1in]{geometry}
\usepackage{float}
\usepackage{graphicx}
\usepackage{tabularx}
\usepackage{mathtools}
\usepackage{tikz}
\usetikzlibrary{positioning}

\date{
	\vspace{-.75em}
	{\large \today}
}

\title{
	\vspace*{-1.5em}
	{\Huge Lecture 19} \\
	\vspace{-1em}
	\rule{\textwidth/4}{.01em} \\
	\vspace{-.25em}
	{\Large Formal Modelling}
	\vspace{-.75em}
}

\author{
	{\large Matthew Farah}
}

\hypersetup{
	colorlinks=true,
	linkcolor=blue,
	filecolor=magenta,
	urlcolor=blue,
	citecolor=blue,
	anchorcolor=blue
}

\fancyhead{}
\fancyhead[L]{\today}
\fancyhead[R]{Lecture 19 - Formal Modelling}

\setlength{\parindent}{0cm}

\lstset{
	basicstyle=\ttfamily\small,
	frame=single,
	tabsize=4,
	breaklines=true,
	numbers=left,
	numberstyle=\tiny,
	numbersep=8pt,
	xleftmargin=1.5em,
	framexleftmargin=1.5em
}

\newcounter{example}
\renewcommand{\theexample}{\arabic{example}}

\newenvironment{example}[1][]{
	\refstepcounter{example}
	\begin{adjustwidth}{1.5em}{}
	\noindent\textbf{\ifx&#1&Example \theexample:\else Example \theexample \ (#1):\fi}
	\ignorespaces
}{
	\vspace{.5em}
	\end{adjustwidth}
}

\newcounter{subexample}
\renewcommand{\thesubexample}{\alph{subexample}}

\newenvironment{subexample}{
	\setcounter{step}{0}
	\refstepcounter{subexample}
	\vspace{1em}\par\noindent
	\begin{adjustwidth}{1.5em}{}
	\noindent\textbf{(\thesubexample)}
	\ignorespaces
}{
	\par
	\vspace{1em}
	\end{adjustwidth}
}

\newenvironment{solution}{
	\vspace{.5em}
	\par
	\begin{adjustwidth}{1.5em}{}
	\textbf{{Solution:}}
	\par
	\vspace{.5em}
}{
	\vspace{1em}
	\end{adjustwidth}
}

\newcounter{step}
\renewcommand{\thestep}{\Roman{step}}

\newenvironment{step}{
	\refstepcounter{step}
	\vspace{1em}\par\noindent
	\ignorespaces\noindent\textbf{(\thestep)}
	\ignorespaces
}{
	\par
	\vspace{1em}
}

\sectionfont{\fontsize{12}{12}\bfseries}
\subsectionfont{\fontsize{10}{12}\normalfont}
\subsubsectionfont{\fontsize{10}{12}\normalfont}

\begin{document}

\maketitle
\pagestyle{fancy}

\fi

%TODO: Add examples, if I have time
\begin{itemize}
	\item Building models of system helps in the testing process.
	\begin{itemize}
		\item Helps abstract away details unnecessary to the testing process.
		\item Allows for the system to be evaluated without an actual implementation.
	\end{itemize}
	\item Models can be:
	\begin{enumerate}
		\item Formal.
		\begin{itemize}
			\item Built around the system's formal requirements.
		\end{itemize}
		\item Informal.
		\begin{itemize}
			\item A rough draft of the system.
		\end{itemize}
	\end{enumerate}
	\item There are many different ways to model the same system.
	\begin{itemize}
		\item Ex. Petri-nets, state-transition systems, etc.
	\end{itemize}
	\item There is no single `best' model for a system.
	\begin{itemize}
		\item The best model for a system is context-dependent.
	\end{itemize}
	\item Model checking a system follows a two-step process:
	\begin{enumerate}
		\item Building a formal model of the system.
		\begin{itemize}
			\item Using the system's formal requirements.
		\end{itemize}
		\item Evaluating the constructed model.
	\end{enumerate}
	\item Model checking is an undecidable problem.
	\begin{itemize}
		\item There doesn't exist an algorithm which can always determine whether a given model of a system always behaves correctly or not.
	\end{itemize}
	\item All languages are composed of at least two things:
	\begin{enumerate}
		\item A concrete syntax.
		\begin{itemize}
			\item The language's notation.
		\end{itemize}
		\item An abstract syntax.
		\begin{itemize}
			\item A set of rules governing how the language's notation can be structured.
		\end{itemize}
	\end{enumerate}
	\item Formalisms are languages with semantics.
	\begin{itemize}
		\item Semantics assign meaning to notation.
	\end{itemize}
	\item Models are expressed using formalisms to prevent ambiguity.
	\item Some notable formal models include:
	\begin{itemize}
		\item Kipke structures.
		\begin{itemize}
			\item TLDR; Basically just finite state machines.
			\item Don't think we really need to know these.
		\end{itemize}
		\item Timed automata.
		\begin{itemize}
			\item TLDR; Kipke structures that additionally have:
			\begin{itemize}
				\item An infinite state-space.
				\item Timed semantics (behaviour depends on transitions and the elapsed time).
			\end{itemize}
			\item Also don't think we really need to know these.
		\end{itemize}
		\item Petri-nets.
		\begin{itemize}
			\item TLDR; See SFWRENG-3BB4.
			\item We do need to know these.
		\end{itemize}
	\end{itemize}
	\item Petri-nets are used to model:
	\begin{itemize}
		\item Concurrency.
		\item Causality.
		\item Synchronization.
		\item Resource constraints.
	\end{itemize}
	\item Petri-nets verify system properties such as:
	\begin{itemize}
		\item Reachability.
		\item The existence of deadlocks.
		\item Boundedness.
		\item Liveliness
		\begin{itemize}
			\item \emph{Aside: The slides call it liveness but I'm using liveliness because liveness sounds so dumb to me.}
		\end{itemize}
	\end{itemize}
\end{itemize}

\ifdefined\iscombined
	\newpage
\else
	\end{document}
\fi
